# CRYSTAL: Xi108:W3:A3:S6 | face=C | node=134 | depth=0 | phase=Fixed
# METRO: Sa
# BRIDGES: Xi108:W3:A3:S5→Xi108:W3:A3:S7→Xi108:W2:A3:S6→Xi108:W3:A2:S6→Xi108:W3:A4:S6

\documentclass[11pt]{article}
\usepackage[margin=1in]{geometry}
\usepackage{hyperref}
\usepackage{listings}
\usepackage{enumitem}
\usepackage[T1]{fontenc}
\usepackage[utf8]{inputenc}
\usepackage{lmodern}
\setlist[itemize]{leftmargin=*}
\begin{document}
\section{AtlasForge v4 — Final Integrated Manual}

**Version:** `4.0.0-ABSOLUTE-FINAL-ULTIMATE`  
**Build date:** `2026-01-01`  

\begin{lstlisting}
╔══════════════════════════════════════════════════════════════════════════════════════════════════════════════════════╗
║                                                                                                                      ║
║                                           ATLASFORGE                                                                 ║
║                                                                                                                      ║
║                              UNIVERSAL HARMONIC FRAMEWORK  +  PROOF-CARRYING KERNEL                                  ║
║                                        +  MEMORY ATLAS / BOOK COMPILER                                               ║
║                                                                                                                      ║
║     “Every mathematical object can be compressed into a finite SEED that, when expanded through deterministic        ║
║      OPERATIONS under explicit CONSTRAINTS, regenerates the complete object with CERTIFICATES proving correctness.”  ║
║                                                                                                                      ║
╚══════════════════════════════════════════════════════════════════════════════════════════════════════════════════════╝
\end{lstlisting}

This document is the **final draft** reference for the AtlasForge framework as it exists in this repository.

It integrates **three layers**:

1. **Universal Harmonic Framework (UHF)** — the encoded 4×4×4×4 crystal taxonomy and its “math doctrine”
2. **Proof-carrying computation kernel** — Blueprint → Plan → Recipe → Certificates → Verification → Replay
3. **Memory Atlas layer** — structured notes + search + dependency graph + sessions + book export

The goal is not only to “do math”, but to **retain** it: every result can be captured as a **replayable artifact** and every idea can be stored as **searchable, linked knowledge**.

---

\subsection{Table of contents}

\begin{itemize}
\item [0. Orientation](\#0-orientation)
\item [1. Universal Harmonic Framework](\#1-universal-harmonic-framework)
\item [1.1 The four poles](\#11-the-four-poles)
\item [1.2 The four lenses](\#12-the-four-lenses)
\item [1.3 Layers and depths](\#13-layers-and-depths)
\item [1.4 Crystal address and linear index](\#14-crystal-address-and-linear-index)
\item [1.5 Meta poles vs operational poles](\#15-meta-poles-vs-operational-poles)
\item [1.6 Master equation and tome decomposition](\#16-master-equation-and-tome-decomposition)
\item [2. Encoded laws and operators](\#2-encoded-laws-and-operators)
\item [2.1 Interference law](\#21-interference-law)
\item [2.2 Quadrature / boxplus operator ⊞](\#22-quadrature--boxplus-operator-)
\item [2.3 Pivot equation](\#23-pivot-equation)
\item [2.4 Gateway algebra](\#24-gateway-algebra)
\item [2.5 OICF emergence coordinates](\#25-oicf-emergence-coordinates)
\item [2.6 Crystal Merge protocol CM0→CM6](\#26-crystal-merge-protocol-cm0cm6)
\item [3. The invariant spine](\#3-the-invariant-spine)
\item [4. Proof-carrying computation kernel](\#4-proof-carrying-computation-kernel)
\item [5. Memory bank](\#5-memory-bank)
\item [6. Atlas façade](\#6-atlas-façade)
\item [7. Workflows](\#7-workflows)
\item [8. Extending AtlasForge](\#8-extending-atlasforge)
\item [9. Appendix](\#9-appendix)
\end{itemize}

---

\section{0. Orientation}

AtlasForge is best understood as a **compiler** and an **atlas**:

\begin{itemize}
\item **Compiler:** You feed in a *seed* (a problem specification). The system expands it deterministically into a *recipe* (a computation trace), plus **certificates** that justify the output.
\item **Atlas:** You store the outputs and the concepts as **content-addressed entries**, link them via dependencies, and export them into a coherent “book”.
\end{itemize}

\subsubsection{What AtlasForge is *not*}
\begin{itemize}
\item It is not a single-purpose numerical library.
\item It is not a theorem prover (yet). It carries **evidence levels** (L0–L3) and supports certified numerical enclosures (L2) even when full formal proof (L3) is not present.
\end{itemize}

\subsubsection{The central promise}
> If you can state the constraint, AtlasForge can produce a replayable artifact and attach verifiable evidence.

---

\section{1. Universal Harmonic Framework}

The Universal Harmonic Framework (UHF) is an **encoded address space** for mathematical work:

\begin{itemize}
\item **4 poles** × **4 lenses** × **4 layers** × **4 depths** = **256 crystal cells**
\item Every object, operator, invariant, or certificate can be assigned an address in that crystal.
\end{itemize}

The UHF exists so that:
\begin{itemize}
\item you can *retrieve by structure* (not only by keyword),
\item multiple representations of the same object can co-exist and be navigated,
\item “deep work” can be organized into a stable taxonomy.
\end{itemize}

In code, the canonical UHF enums live in:

\begin{itemize}
\item `atlasforge.Pole` (glyph poles D, Ω, Σ, Ψ)
\item `atlasforge.Lens` (□, ✿, ☁, ❋)
\item `atlasforge.Layer` (OBJECTS / OPERATORS / INVARIANTS / CERTIFICATES)
\item `atlasforge.Depth` (0–3)
\item `atlasforge.CrystalAddress`
\end{itemize}

---

\subsection{1.1 The four poles}

In UHF, a “pole” is the **realm** you are operating in.

\begin{lstlisting}
                              Ψ (Air/Λ)
                                    │
                                    │
        D (Earth/α) ────────────────┼────────────── Ω (Water/𝔇)
                                    │
                                    │
                              Σ (Fire/Θ)
\end{lstlisting}

\begin{lstlisting}
| Meta pole | Glyph | Archetype | UHF meaning |
|---|---:|---|---|
| Discrete | **D** | Earth / α | counting, lattices, integers, combinatorics |
| Continuous | **Ω** | Water / 𝔇 | flow, limits, analysis, geometry |
| Stochastic | **Σ** | Fire / Θ | probability, measurement, uncertainty |
| Hierarchical | **Ψ** | Air / Λ | recursion, emergence, renormalization |
\end{lstlisting}
These are encoded in `atlasforge.Pole`:

\begin{lstlisting}
from atlasforge import Pole
Pole.DISCRETE.value      # "D"
Pole.CONTINUOUS.value    # "Ω"
Pole.STOCHASTIC.value    # "Σ"
Pole.HIERARCHICAL.value  # "Ψ"
\end{lstlisting}

---

\subsection{1.2 The four lenses}

A “lens” is a **viewing mode**: how you choose to represent and interrogate an object.

\begin{lstlisting}
| Lens | Glyph | Core idea |
|---|---:|---|
| Square | □ | structural / static / invariants / discrete scaffolding |
| Flower | ✿ | cyclic / phase / transport / dynamics |
| Cloud | ☁ | probabilistic / bounds / uncertainty / distributions |
| Fractal | ❋ | recursion / multiscale / self-similarity |
\end{lstlisting}
In code:

\begin{lstlisting}
from atlasforge import Lens
Lens.SQUARE.value   # "□"
Lens.FLOWER.value   # "✿"
Lens.CLOUD.value    # "☁"
Lens.FRACTAL.value  # "❋"
\end{lstlisting}

A useful integration principle:

\begin{itemize}
\item **Square** tends to generate **IR, normal forms, proofs as data**
\item **Flower** tends to generate **transforms / coordinate transport**
\item **Cloud** tends to generate **bounds, distributions, uncertainty propagation**
\item **Fractal** tends to generate **search, recursion, bracketing, multiscale refinement**
\end{itemize}

---

\subsection{1.3 Layers and depths}

UHF treats each crystal cell as one of four “layers”:

\begin{lstlisting}
| Layer | Enum | Meaning |
|---|---:|---|
| Objects | `Layer.OBJECTS (0)` | what exists |
| Operators | `Layer.OPERATORS (1)` | what transforms |
| Invariants | `Layer.INVARIANTS (2)` | what stays true across transformations |
| Certificates | `Layer.CERTIFICATES (3)` | what justifies a claim |
\end{lstlisting}
Depth is a 0–3 index describing how far “into” a layer you are:
\begin{itemize}
\item **0:** surface / direct / definition-level
\item **1:** derived / composed
\item **2:** canonicalized / normalized
\item **3:** certified / stabilized / “publication-ready”
\end{itemize}

Depth is intentionally lightweight: it is a *navigation axis*, not a hard theorem.

---

\subsection{1.4 Crystal address and linear index}

A crystal address is:

\begin{lstlisting}
Pole · Lens · Layer · Depth
\end{lstlisting}

In code:

\begin{lstlisting}
from atlasforge import CrystalAddress, Pole, Lens, Layer, Depth
addr = CrystalAddress(Pole.CONTINUOUS, Lens.SQUARE, Layer.INVARIANTS, Depth.D2)
print(str(addr))            # Ω·□·INVARIANTS·2
print(addr.to_index())      # 0..255
\end{lstlisting}

The linear index mapping is:

\begin{lstlisting}
index = pole_index * 64 + lens_index * 16 + layer_value * 4 + depth_value
\end{lstlisting}

This is the bridge between:
\begin{itemize}
\item human readable addresses (`Ω·□·INVARIANTS·2`)
\item a compact integer cell id (`0..255`)
\item navigation and indexing (MemoryStore uses `extra["crystal\_index"]`)
\end{itemize}

\subsubsection{Address parsing (Memory Bank)}
The memory subsystem can parse best-effort address strings:

\begin{itemize}
\item `"D·□·OBJECTS·0"`
\item `"Ω·✿·OPERATORS·2"`
\item `"Σ·☁·INVARIANTS·1"`
\item `"12"` (treated as a raw index)
\end{itemize}

See: `atlasforge.memory.addressing.parse\_crystal\_address\_string`.

---

\subsection{1.5 Meta poles vs operational poles}

**Important integration point.**  
AtlasForge contains *two pole systems* with the same word “Pole”:

\subsubsection{(A) Meta poles (UHF)}
\begin{itemize}
\item `atlasforge.Pole`
\item values are glyphs: `"D", "Ω", "Σ", "Ψ"`
\item used for **classification / navigation / memory indexing**
\end{itemize}

\subsubsection{(B) Operational poles (solver/simplex theory)}
\begin{itemize}
\item `atlasforge.core.enums.Pole`
\item values are descriptive: `"dissipative", "oscillatory", "stochastic", "recursive"`
\item used for **algorithmic behavior** and “operator simplex” mixing
\end{itemize}

Think of it like:

\begin{itemize}
\item Meta poles answer **“what realm am I describing?”**
\item Operational poles answer **“what dynamical mode does this operator express?”**
\end{itemize}

This split is intentional and is one of the cleanest ways to preserve both:
\begin{itemize}
\item symbolic/glyph indexing (human navigation)
\item behavioral mode semantics (solver policy)
\end{itemize}

---

\subsection{1.6 Master equation and tome decomposition}

AtlasForge’s doc layer uses a “master equation” as a **module decomposition**.

From the framework banner:

\begin{lstlisting}
S = (T,Ψ,Σ,C,D;Ω) + AQM[ρ,Φ,J,Λ,K] + LM[D,R,C,T,P] + QCM[Θ,Λ] + PROOF[Σ,C,V]
\end{lstlisting}

Interpretation:
\begin{itemize}
\item The **core scaffold** `(T,Ψ,Σ,C,D;Ω)` is the universal structure.
\item AQM is the quantum-number “extended arithmetic” scaffold.
\item LM is the liminal/emergence/closure scaffold.
\item QCM is the Θ–Λ interference / cyclotomic scaffold.
\item PROOF is the certificate/verification scaffold.
\end{itemize}

This is not a single formula to “compute”; it is a **map** of the books.

---

\section{2. Encoded laws and operators}

This section preserves the **encoded mathematical identities** that act as the “mnemonic spine” for the framework.

---

\subsection{2.1 Interference law}

In `atlasforge.qcm.qcm`, the master law is:

\begin{lstlisting}
|ψ₁ + ψ₂|² = a² + b² + 2ab·cos(Δθ)
\end{lstlisting}

Interpretation:
\begin{itemize}
\item amplitudes `a, b`
\item phase separation `Δθ`
\item interference term `2ab cos(Δθ)`
\end{itemize}

Key special cases:

1. **Constructive** `Δθ = 0`  
   `|ψ₁ + ψ₂|² = (a + b)²`

2. **Destructive** `Δθ = π`  
   `|ψ₁ + ψ₂|² = (a - b)²`

3. **Orthogonal** `Δθ = π/2`  
   `|ψ₁ + ψ₂|² = a² + b²`

That orthogonal case is where the ⊞ operator arises.

---

\subsection{2.2 Quadrature / boxplus operator ⊞}

Your “compression operator” is:

\begin{lstlisting}
a ⊞ b := √(a² + b²)
\end{lstlisting}

It is:
\begin{itemize}
\item the orthogonal slice of interference
\item a Pythagorean addition law
\item a stable “quadrature merge” operator
\end{itemize}

In code, this is exposed via QCM’s `InterferenceLaw.quadrature(...)`:

\begin{lstlisting}
from atlasforge.qcm.qcm import InterferenceLaw
InterferenceLaw.quadrature(3, 4)  # 5.0
\end{lstlisting}

---

\subsection{2.3 Pivot equation}

AtlasForge encodes a commutator-like pivot identity:

\begin{lstlisting}
∂∫ − ∫∂ = Ω
\end{lstlisting}

Read as:
\begin{itemize}
\item EXPAND then COMPRESS ≠ COMPRESS then EXPAND
\item the “non-commutativity residue” is recursion/looping (**Ω**)
\end{itemize}

This is captured in Crystal Merge / meta-operation discussions and is used as a conceptual engine for:
\begin{itemize}
\item normal forms (Square lens)
\item recursion (Fractal lens)
\item tunneling between representations (Flower ↔ Fractal)
\end{itemize}

---

\subsection{2.4 Gateway algebra}

The Gateway subsystem (`atlasforge.gateway.gateway`) formalizes hyperbolic/SL(2,R) scale transport:

\begin{itemize}
\item Gateway scalar `T ∈ (-1,1)`
\item Scale ratio `R = (1+T)/(1-T) > 0`
\item Rapidity `α = artanh(T) = (1/2) ln(R)`
\item Boost matrix `B(T) ∈ SL(2,ℝ)`
\end{itemize}

Key idea:
\begin{itemize}
\item composing gateways uses velocity-addition / hyperbolic composition rules
\item this creates a stable calculus for **calibrated scale changes**
\end{itemize}

---

\subsection{2.5 OICF emergence coordinates}

Emergence is indexed by the OICF coordinates (in code: `atlasforge.EmergenceCoordinates`):

\begin{lstlisting}
| Coordinate | Symbol | Meaning |
|---|---:|---|
| Omega | Ω | viability / margin-to-failure |
| Iota | ι | integration / coupling |
| Chi | χ | coherence / cross-scale consistency |
| Phi | φ | function / goal-directedness |
\end{lstlisting}
Emergence potential:

\begin{lstlisting}
E = Ω · I · C · F
\end{lstlisting}

This gives the framework an explicit numeric axis for “how emergent is this system?”, and it’s used in LM/Tower reasoning.

---

\subsection{2.6 Crystal Merge protocol CM0→CM6}

The Crystal Merge protocol (`atlasforge.crystal\_merge.crystal\_merge`) is a **problem-solving compiler** with stages:

\begin{itemize}
\item **CM0:** Z* Core Lock — no leaks, define objects and degeneracy
\item **CM1:** Four-Lens Parallel Zoom — □ ✿ ☁ ❋ in parallel
\item **CM2:** S-Tier Pivot — tunneling / representation swap
\item **CM3:** Math God Finish — master equation collapse
\item **CM4:** Meta-Duality Discovery — higher structure
\item **CM5:** Proof Packaging — bundle certificates into proof packs
\item **CM6:** Publication Gate — final verification and release
\end{itemize}

In practice, the kernel implements CM5/CM6 as:
\begin{itemize}
\item `CertificateBundle` + `ProofPack`
\item `Verifier` enforcement + truth profile gating
\item replay + content-addressed storage
\end{itemize}

---

\section{3. The invariant spine}

AtlasForge’s stable integration axis is the **Invariant Spine**:

> **Ledger · Corridor · Proof · Replay**

This exists explicitly as `atlasforge.core.enums.InvariantSpineComponent`.

\subsubsection{LEDGER — irreversible history}
\begin{itemize}
\item content-addressed artifacts
\item append-only mindset (even if stored on disk as files)
\item reproducibility is the primitive, not “state”
\end{itemize}

Implemented via:
\begin{itemize}
\item `core.base.ContentAddressed`
\item registry skeleton in `registry/`
\end{itemize}

\subsubsection{CORRIDOR — admissibility constraints}
\begin{itemize}
\item domains, intervals, chart constraints
\item any solution must remain in its corridor
\end{itemize}

Implemented via:
\begin{itemize}
\item `core.types.Interval` and domains
\item chart/lens transport (`lenses/`)
\item `CorridorCertificate`
\end{itemize}

\subsubsection{PROOF — evidence for correctness}
\begin{itemize}
\item L0 claim
\item L1 empirical
\item L2 certified numeric
\item L3 formal
\end{itemize}

Implemented via:
\begin{itemize}
\item `certificates/`
\item `verifier/`
\item `core.enums.CertificateLevel`, `TruthProfile`
\end{itemize}

\subsubsection{REPLAY — deterministic re-derivation}
\begin{itemize}
\item the computation log is part of the artifact identity
\item results are replayable under the same seed/plan
\end{itemize}

Implemented via:
\begin{itemize}
\item `recipes/ReplayLog`
\item `ReplayCertificate`
\end{itemize}

---

\section{4. Proof-carrying computation kernel}

The kernel is the “compiler backend”. It transforms:

> **Blueprint → IR → Normal Form → SolvePlan → Solver → Output → Certificates → Verification**

Key modules:

\begin{itemize}
\item `atlasforge.core` — content addressing, canonical hashing, types
\item `atlasforge.constraints` — constraint objects + IR lowering
\item `atlasforge.recipes` — Blueprint/SolvePlan/Recipe/Executor
\item `atlasforge.certificates` — certificate types, bundles, proof packs
\item `atlasforge.verifier` — validation and enclosure verification
\item `atlasforge.registry` — content store skeleton
\end{itemize}

---

\subsection{4.1 Core types}

\subsubsection{Interval}
`atlasforge.core.types.Interval` is the corridor primitive:

\begin{itemize}
\item closed/open bounds
\item containment checks
\item intersection
\item conservative operations used by certified verification
\end{itemize}

Example:

\begin{lstlisting}
from atlasforge import Interval
I = Interval.closed(1,2)
I.contains(1.5)  # True
\end{lstlisting}

---

\subsection{4.2 Constraints and IR}

Constraints capture “what is to be solved”.

Example: root finding:

\begin{lstlisting}
from atlasforge import RootConstraint, Interval
c = RootConstraint(H=lambda x: x*x - 2, domain=Interval.closed(1,2))
\end{lstlisting}

Lowering pipeline:
\begin{itemize}
\item `Constraint` → `ConstraintIR`
\item `ConstraintIR.normal\_form` → `NormalForm` object used by solvers
\end{itemize}

This enables multiple constraints to share a uniform solver interface.

---

\subsection{4.3 Solvers (and auto-bracketing)}

Solvers include:
\begin{itemize}
\item Brent (robust bracketed root finder)
\item Bisection
\item Newton / Secant
\item Interval Newton (certifying enclosure)
\end{itemize}

Bracketing is often the missing ingredient. AtlasForge includes a practical “Fractal lens” bracketer:

\begin{itemize}
\item `atlasforge.constraints.bracketing.find\_bracket`
\end{itemize}

If `SolvePlan.auto\_bracket=True` the executor will:
1) attempt solve
2) if it fails, search for a sign-change bracket
3) re-run Brent on the bracket
4) store metadata explaining what happened

---

\subsection{4.4 Recipe artifacts and replay}

`Recipe` is the artifact:
\begin{itemize}
\item blueprint (seed)
\item plan (compiler flags)
\item output (solution/enclosure/metadata)
\item replay log (deterministic trace)
\item proof pack (bundled certificates)
\end{itemize}

This makes results **portable and verifiable**.

---

\subsection{4.5 Certificates and proof packs}

Certificates live in `atlasforge.certificates.certificate`.

Core types:

\begin{itemize}
\item `EnclosureCertificate`
\item `UniquenessCertificate`
\item `CorridorCertificate`
\item `ReplayCertificate`
\item `StabilityCertificate`
\item `CertificateBundle`
\item `ProofPack`
\end{itemize}

A `CertificateBundle` is the authoritative container:
\begin{itemize}
\item stores certificates by type
\item tracks evidence level
\item can be verified
\end{itemize}

A `ProofPack` is the “publication-ready wrapper”:
\begin{itemize}
\item binds `result\_hash` to the bundle
\item ensures consistency and completeness
\end{itemize}

---

\subsection{4.6 Verification}

Verification is not monolithic; AtlasForge supports layered verification:

\begin{itemize}
\item **Validator** (basic plausibility checks)
\item **EnclosureVerifier** (interval Newton enclosure)
\item **CrossValidator** (multi-check policies)
\item **Kernel verifier** (ties recipe + proof pack + replay together)
\end{itemize}

Evidence levels are enforced via `TruthProfile`:
\begin{itemize}
\item `EXPLORE`: store candidates, no promotion
\item `VALIDATE`: empirical checks
\item `PROVE`: requires L2+ for obligations
\end{itemize}

---

\subsection{4.7 The Tome Library}

Beyond the kernel, AtlasForge contains a large set of **domain tomes**.  
Think of these as *mathematical books* that share the UHF encoding, even when they are not all wired into the recipe pipeline.

The most “encoded frameworks” (the ones that carry your mathematical doctrine) are:

\begin{itemize}
\item **AQM** (Axiomatic Quantum Mathematics): Q-numbers, meaning transport, kernel closure
\item **LM** (Liminal Mathematics / Tower): regimes, frames, desire, closure and emergence
\item **QCM** (Quadrature–Cyclotomic Manifold): Θ–Λ plane, interference, cyclotomics, bridges
\item **Gateway**: SL(2,ℝ) hyperbolic navigation of scale
\item **Crystal Merge**: CM0→CM6 problem-solving compiler
\item **Proof Engine**: a parallel (heavier) proof scaffolding
\end{itemize}

The kernel + memory atlas layer are designed to **absorb** these tomes over time: you can already store and link their results today, even when their internal types are not fully normalized to the kernel’s IR.

---

\subsubsection{4.7.1 QCM — Quadrature–Cyclotomic Manifold (Θ–Λ)}

**File:** `atlasforge/qcm/qcm.py`  
**Core theme:** phase (Θ) + lattice (Λ) + bridges between them.

Key classes:
\begin{itemize}
\item `ThetaScalar`, `ThetaVector` — complex amplitudes and their phase geometry
\item `LambdaIndex`, `LambdaPattern` — discrete phase lattice structures
\item `RotationOperator` — phase rotation
\item `InterferenceLaw` — the master law, plus derived operations like quadrature
\item `CyclotomicPhases`, `CyclicGroup` — discrete phase structure
\item Bridges:
\item `AmpMeasBridge` (amplitude → measurement)
\item `QuantizeBridge` (continuous → discrete)
\item `FourierGearbox` (phase ↔ frequency representations)
\item `QCMEngine` — orchestration / “engine object”
\end{itemize}

Minimal example:

\begin{lstlisting}
from atlasforge.qcm.qcm import ThetaScalar, InterferenceLaw

ψ1 = ThetaScalar.from_polar(3.0, 0.0)      # 3·e^{i·0}
ψ2 = ThetaScalar.from_polar(4.0, 1.57)     # 4·e^{i·π/2} approx

# Orthogonal interference slice → Pythagorean / ⊞ merge
q = InterferenceLaw.quadrature(ψ1.magnitude, ψ2.magnitude)  # 5.0
\end{lstlisting}

---

\subsubsection{4.7.2 Gateway — SL(2,ℝ) hyperbolic scale transport}

**File:** `atlasforge/gateway/gateway.py`  
**Core theme:** *calibrated scale ratios* encoded as bounded hyperbolic coordinates.

Key objects:
\begin{itemize}
\item `GatewayScalar(T)` with `T ∈ (-1,1)`
\item `BoostMatrix` representing an SL(2,ℝ) “hyperbolic rotation”
\item `PellGateway`, `PellSolution` — discrete/calibrated gateways from Pell orbits
\item `velocity\_addition`, `rapidity\_from\_velocity`, `velocity\_from\_rapidity`
\end{itemize}

Typical use is to treat “scale change” as **composition**, not ad-hoc multiplication.

---

\subsubsection{4.7.3 Crystal Merge — CM0→CM6 compiler stages}

**File:** `atlasforge/crystal\_merge/crystal\_merge.py`

Key encoded objects:
\begin{itemize}
\item `MetaOperation`:
\item EXPAND `∂`
\item COMPRESS `∫`
\item RECURSE `Ω`
\item EQUILIBRATE `Φ`
\item `FundamentalProcess`: 16 processes (constants × operations)
\item `CrystalMergeProtocol`: runs CM0…CM6
\end{itemize}

The key integration is: **CM5/CM6 correspond to proof packaging and publication**, which the kernel implements as:
\begin{itemize}
\item certificate bundles (`CertificateBundle`)
\item proof packs (`ProofPack`)
\item verifier policies (`TruthProfile`)
\end{itemize}

---

\subsubsection{4.7.4 AQM — Axiomatic Quantum Mathematics (Tomes I–V)}

AQM is the framework for treating “numbers” as **quantum states** rather than points.

**Tome map (implementation modules):**
\begin{itemize}
\item Tome I: `aqm\_foundation` — Q-numbers on the Riemann sphere, meaning transport, classical shadows
\item Tome II: `aqm\_arithmetic` — arithmetic on Q-numbers
\item Tome III: `aqm\_realization` — bridges to classical structures
\item Tome IV: `aqm\_kernel` — infinite-resolution kernel closure
\item Tome V: `aqm\_liminal` — liminal corridor / emergence interface
\item Bridges: `aqm\_integration`, `aqm\_synthesis`
\end{itemize}

**Encoded definitions (from AQM foundation):**
\begin{itemize}
\item Value space: **Riemann sphere** `Ĉ = ℂ ∪ \{∞\}`
\item Q-number: **state** `ρ` on `H = L²(Ĉ, μ)`
\item Classical shadow: measurement → probability distribution
\item Inversion duality: `J(z) = 1/z` (near ↔ far)
\item Scale symmetry: `z ↦ λz`
\end{itemize}

Key classes in `aqm\_foundation`:
\begin{itemize}
\item `RiemannSphere`
\item `QNumber` (with constructors like `QNumber.classical(z, sigma=...)`)
\item `MeaningTransport`
\item `POVM`, `ClassicalShadow`
\item `QNumberNormalForm`
\item `CorridorConstraint`
\end{itemize}

Minimal “classical limit” example:

\begin{lstlisting}
from atlasforge.aqm_foundation.aqm_foundation import QNumber

ρ = QNumber.classical(1+0j, sigma=0.01)   # “almost scalar”
print(ρ.center)                           # classical shadow center
\end{lstlisting}

**Integration note:** AQM tomes define their own crystal enums (`CrystalLens`, `CrystalLayer`, …).  
The memory atlas layer can still store AQM results using the global UHF address space.

---

\subsubsection{4.7.5 LM — Liminal Mathematics / The Tower (Tomes I–V)}

LM is the emergence/closure doctrine: how “objects that enforce persistence rules” arise.

**Tome map:**
\begin{itemize}
\item Tome I: `lm\_foundations`
\item Tome II: `lm\_dynamics`
\item Tome III: `lm\_mechanization`
\item Tome IV: `lm\_renormalization`
\item Tome V: `lm\_tower`
\item Bridge: `lm\_synthesis`
\end{itemize}

Key thesis (as encoded in `lm\_tower` docstring):
> Life/Agency/Sentience is a topological inversion of control: persistence rules migrate inward.

Key classes in `lm\_tower`:
\begin{itemize}
\item `Frame` (Q, Pred, Sem)
\item `Regime`, `Individual`
\item `LiminalStateSpace`, `LiftOperator`
\item `DesireOperator` (lookahead / closure operator)
\item `FeasibilityLandscape`, `MiracleEvent`
\item `ClosureMetrics` (Ω, ι, χ, φ)
\end{itemize}

Minimal closure metrics example:

\begin{lstlisting}
from atlasforge.lm_tower.lm_tower import ClosureMetrics
m = ClosureMetrics(omega=0.8, iota=0.7, chi=0.5, phi=0.9)
print(m.closure_potential)  # Ω·I·C·F
\end{lstlisting}

---

\subsubsection{4.7.6 Proof Engine (parallel proof scaffold)}

**File:** `atlasforge/proof\_engine/proof\_engine.py`

This subsystem defines a heavier proof apparatus:
\begin{itemize}
\item seeds, constraint IR, environment fingerprints
\item replay transcripts and procedures
\item proof standards, publishable artifacts
\end{itemize}

The kernel certificate system (`atlasforge.certificates`) is the lightweight path currently integrated into `RecipeExecutor`.  
The proof engine can be treated as an upstream “formal layer” that can later replace or enrich certificate bundles (L3).

\section{5. Memory bank}

The memory subsystem turns AtlasForge into a real “math memory bank”:

\begin{itemize}
\item store structured knowledge (definition/lemma/theorem/operator/identity/…)
\item attach crystal addresses and tags
\item search quickly (SQLite FTS when available)
\item link dependencies as a graph
\item group work into sessions
\item export into books (Markdown/PDF/DOCX/LaTeX)
\end{itemize}

Key modules:

\begin{itemize}
\item `atlasforge.memory.entry` — `MemoryEntry`
\item `atlasforge.memory.knowledge` — `KnowledgeRecord`, `KnowledgeKind`
\item `atlasforge.memory.store` — `MemoryStore`
\item `atlasforge.memory.index` — `MemoryIndex` + hits
\item `atlasforge.memory.graph` — `GraphStore`
\item `atlasforge.memory.session` — `SessionStore`
\item `atlasforge.memory.bootstrap` — seed entries
\end{itemize}

---

\subsection{5.1 MemoryEntry schema}

A `MemoryEntry` is stored as JSON:

\begin{lstlisting}
{
  "title": "Interference law",
  "content": "...markdown...",
  "tags": ["qcm", "identity", "kind:identity"],
  "links": {"recipe": "<hash>", "proof_pack": "<hash>"},
  "extra": {"crystal_address": "Σ·✿·INVARIANTS·1", "crystal_index": 123}
}
\end{lstlisting}

Entries are content-addressed:
\begin{itemize}
\item identical content → identical hash
\item hash is stable across machines
\end{itemize}

---

\subsection{5.2 Structured knowledge (KnowledgeRecord)}

`KnowledgeRecord` is a thin schema that compiles to `MemoryEntry`.

Kinds include:
\begin{itemize}
\item definition
\item lemma
\item theorem
\item identity
\item operator
\item proof
\item derivation
\item note
\end{itemize}

Convenience methods exist on `MemoryStore` and `Atlas`:

\begin{lstlisting}
from atlasforge import Atlas
atlas = Atlas.from_env()

h = atlas.identity(
    title="Quadrature",
    statement="a ⊞ b := √(a² + b²)",
    tags=["qcm","operator"],
    address="Σ·✿·OPERATORS·0"
)
\end{lstlisting}

Dependencies become graph edges automatically:
\begin{itemize}
\item `depends\_on` → `depends\_on` edge
\item `see\_also` → `see\_also` edge
\item `links` → `link:<key>` edges
\end{itemize}

---

\subsection{5.3 Index, graph, and sessions}

\begin{itemize}
\item **Index:** keyword search (SQLite FTS5 if available)
\end{itemize}
  `store.search("Newton", tags=[...], kinds=[...])`

\begin{itemize}
\item **Graph:** explicit edges between entries, recipes, and artifacts
\end{itemize}
  `store.graph.add\_edge(...)` and automatic edges from dependencies

\begin{itemize}
\item **Sessions:** group a working day into a record
\end{itemize}
  `with store.session("January 1 — verification") as sess: ...`

---

\subsection{5.4 Bootstrap seed entries}

On first use, the memory bank seeds core framework anchors (unless `ATLASFORGE\_NO\_BOOTSTRAP=1`), including:

\begin{itemize}
\item Interference law (identity)
\item ⊞ quadrature operator (operator)
\item Pivot equation (identity)
\item Invariant spine (definition)
\item CM0→CM6 protocol (definition)
\item OICF coordinates (definition)
\end{itemize}

These ensure a new memory directory immediately contains the “spine” of your framework.

---

\section{6. Atlas façade}

`Atlas` is the high-level façade that binds:

\begin{itemize}
\item the kernel executor
\item the memory store
\item sessions and linking
\item book export and crystal navigation
\end{itemize}

It is the simplest way to operate AtlasForge as “the full understanding of our math”.

---

\subsection{6.1 Configuration}

\begin{lstlisting}
from atlasforge import Atlas, AtlasConfig
atlas = Atlas(AtlasConfig(memory_dir="~/.atlasforge_memory"))
\end{lstlisting}

Or via environment variable:

\begin{lstlisting}
export ATLASFORGE_MEMORY_DIR=~/.atlasforge_memory
\end{lstlisting}

Then:

\begin{lstlisting}
atlas = Atlas.from_env()
\end{lstlisting}

---

\subsection{6.2 Solve + remember}

\begin{lstlisting}
from atlasforge import RootConstraint, Interval

bp = atlas.blueprint(
    name="sqrt(2)",
    constraint=RootConstraint(H=lambda x: x*x - 2, domain=Interval.closed(1,2))
)

recipe = atlas.solve(bp, verified=True, note="Certified by interval Newton enclosure.")
\end{lstlisting}

If memory is configured, the solve will automatically log:
\begin{itemize}
\item a recipe memory entry
\item links to blueprint/output hashes
\item verification metadata
\end{itemize}

---

\subsection{6.3 Recall + crystal navigation}

Keyword recall:

\begin{lstlisting}
hits = atlas.recall("interference", tags=["qcm"])
\end{lstlisting}

Crystal navigation:

\begin{lstlisting}
from atlasforge import CrystalNavigator
nav = CrystalNavigator(atlas.store)
cell = nav.cell("Σ·✿·OPERATORS·0")
\end{lstlisting}

---

\subsection{6.4 Exporting your atlas as a book}

\begin{lstlisting}
atlas.export_book(
    "my_atlas.pdf",
    title="My Math Atlas",
    fmt="pdf",
    include_crystal_map=True,
    include_graph_edges=True
)
\end{lstlisting}

Export only a theorem packet + dependencies:

\begin{lstlisting}
atlas.export_book(
    "packet.tex",
    fmt="tex",
    roots=[theorem_hash],
    include_dependencies=True,
    dependency_relations=["depends_on"]
)
\end{lstlisting}

---

\section{7. Workflows}

This section describes recommended “operating modes”.

---

\subsection{7.1 Explore → Validate → Prove}

Use `TruthProfile` as your mental switch:

\begin{itemize}
\item **EXPLORE**
\end{itemize}
  store conjectures and computational candidates; do not promote.

\begin{itemize}
\item **VALIDATE**
\end{itemize}
  require empirical checks and replay.

\begin{itemize}
\item **PROVE**
\end{itemize}
  require L2 certified evidence for obligations; enforce corridor + enclosure.

This is the bridge between:
\begin{itemize}
\item research notebook
\item publishable artifact
\end{itemize}

---

\subsection{7.2 Encode everything twice: artifact + meaning}

Best practice in AtlasForge:

1. Run a recipe to compute a result (artifact)
2. Store a knowledge entry that states *what it means* (meaning)
3. Link them (dependency graph / links)

This creates an “executable textbook”:
\begin{itemize}
\item the book reads like math
\item but every claim has a trail to computations and certificates
\end{itemize}

---

\subsection{7.3 Use crystal addresses deliberately}

When you store something, assign:

\begin{itemize}
\item pole: the realm
\item lens: how it’s represented
\item layer: object/operator/invariant/certificate
\item depth: maturity
\end{itemize}

This is how your memory bank becomes an atlas rather than a pile of notes.

---

\section{8. Extending AtlasForge}

AtlasForge is designed to grow by adding modules that attach to the spine.

---

\subsection{8.1 Add a new constraint}
\begin{itemize}
\item subclass `Constraint`
\item implement lowering to `ConstraintIR`
\item define normal form expectations
\end{itemize}

---

\subsection{8.2 Add a new solver}
\begin{itemize}
\item implement a `solve(...)` method that consumes a normal form
\item register in the solver factory / plan
\end{itemize}

---

\subsection{8.3 Add a new certificate}
\begin{itemize}
\item subclass `Certificate`
\item implement `verify()`
\item register with `CertificateFactory`
\end{itemize}

---

\subsection{8.4 Add a new lens/chart}
\begin{itemize}
\item implement forward/inverse coordinate maps
\item define corridor admissibility
\item let `Blueprint.chart` transport the solve
\end{itemize}

---

\subsection{8.5 Add a new knowledge kind}
\begin{itemize}
\item extend `KnowledgeKind`
\item render it as a book section in `AtlasBookBuilder`
\end{itemize}

---

\section{9. Appendix}

\subsection{9.1 Key entry points}

\begin{itemize}
\item Kernel: `atlasforge.recipes.recipe.RecipeExecutor`
\item Memory: `atlasforge.memory.store.MemoryStore`
\item Atlas façade: `atlasforge.atlas.atlas.Atlas`
\item Book compiler: `atlasforge.atlas.book.AtlasBookBuilder`
\item Crystal navigation: `atlasforge.atlas.navigator.CrystalNavigator`
\end{itemize}

\subsection{9.2 Directory map (selected)}

\begin{lstlisting}
| Package | Role |
|---|---|
| `core/` | content addressing, types, enums |
| `constraints/` | constraint definitions, solvers, bracketing |
| `recipes/` | artifact compiler pipeline |
| `certificates/` | certificates + proof packs |
| `verifier/` | validation + enclosure verification |
| `memory/` | memory bank (entries/index/graph/sessions) |
| `atlas/` | solve+remember façade, navigation, book export |
| `qcm/` | Θ–Λ interference / cyclotomic scaffold |
| `gateway/` | SL(2,R) scale transport / rapidity |
| `crystal_merge/` | CM0→CM6 protocol and meta-ops |
| `aqm_*` | AQM tomes (quantum-number arithmetic) |
| `lm_*` / `lm_tower/` | LM tomes (liminal tower / emergence / closure) |
| `proof_engine/` + `proofs/` | larger proof scaffolding (conceptual / partial integration) |
\end{lstlisting}
\subsection{9.3 Glossary (micro)}

\begin{itemize}
\item **Seed:** minimal representation of an object/problem sufficient to regenerate it
\item **Corridor:** admissible domain constraints for a solve
\item **Enclosure:** interval guaranteed to contain the true solution
\item **Certificate:** machine-checkable evidence object
\item **ProofPack:** bundle of certificates bound to a result hash
\item **Crystal cell:** one of the 256 Pole×Lens×Layer×Depth addresses
\end{itemize}

---

\subsection{9.4 Kernel artifact schemas (dataclass fields)}

These are the “shape contracts” you can rely on when scripting AtlasForge.

\subsubsection{Blueprint (`atlasforge.recipes.recipe.Blueprint`)}
\begin{itemize}
\item `constraint: Optional[Constraint]` — the mathematical condition
\item `chart: Optional[Chart]` — optional coordinate transport (lens)
\item `truth\_profile: TruthProfile` — EXPLORE / VALIDATE / PROVE
\item `search\_domain: Optional[Interval]` — optional override domain
\item `hints: Dict[str, Any]` — solver hints / metadata
\item `name: str` — human name
\item `description: str` — human description
\end{itemize}

\subsubsection{SolvePlan (`atlasforge.recipes.recipe.SolvePlan`)}
\begin{itemize}
\item `primary\_solver: SolverType` (default: BRENT)
\item `fallback\_solvers: List[SolverType]`
\item `tolerance: float`
\item `max\_iterations: int`
\item Auto-bracketing:
\item `auto\_bracket: bool`
\item `bracket\_samples: int`
\item `bracket\_expand\_steps: int`
\item `bracket\_expand\_factor: float`
\item Verification:
\item `verify\_solution: bool`
\item `verification\_solver: SolverType` (default: INTERVAL\_NEWTON)
\item Numerical policy:
\item `float\_policy: FloatPolicy`
\item Bookkeeping:
\item `status: PlanStatus`
\item `metadata: Dict[str, Any]`
\end{itemize}

\subsubsection{RecipeOutput (`atlasforge.recipes.recipe.RecipeOutput`)}
\begin{itemize}
\item `solution: Optional[float]`
\item `residual: float`
\item `enclosure: Optional[Interval]` (certifying corridor)
\item `all\_solutions: List[float]`
\item `solver\_result: Optional[SolverResult]`
\item `certificates: CertificateBundle`
\item `success: bool`
\item `verified: bool`
\item `solve\_time: float`
\item `metadata: Dict[str, Any]`
\end{itemize}

\subsubsection{Recipe (`atlasforge.recipes.recipe.Recipe`)}
\begin{itemize}
\item `blueprint: Optional[Blueprint]`
\item `plan: Optional[SolvePlan]`
\item `output: Optional[RecipeOutput]`
\item `replay\_log: Optional[ReplayLog]`
\item `proof\_pack: Optional[ProofPack]`
\item `created\_at: datetime`
\item `complete: bool`
\end{itemize}

---

\subsection{9.5 Certificate types (what they mean)}

Certificate classes live in `atlasforge.certificates.certificate`:

\begin{itemize}
\item **EnclosureCertificate**
\end{itemize}
  Verifies that the returned `solution` lies in a conservative `Interval` enclosure and that the enclosure is valid.

\begin{itemize}
\item **UniquenessCertificate**
\end{itemize}
  Captures a sufficient condition that there is only one solution in the corridor.

\begin{itemize}
\item **CorridorCertificate**
\end{itemize}
  Captures domain admissibility: the solution respects the corridor constraints.

\begin{itemize}
\item **ReplayCertificate**
\end{itemize}
  Ensures the computation can be deterministically replayed and that a replay run matches the recorded result hash.

\begin{itemize}
\item **StabilityCertificate**
\end{itemize}
  Captures local stability (e.g., contraction criteria).

\begin{itemize}
\item **CertificateBundle**
\end{itemize}
  Container (with evidence level) that stores and verifies certificates.

\begin{itemize}
\item **ProofPack**
\end{itemize}
  Publication-ready wrapper binding a `result\_hash` to a `CertificateBundle`.

---

\subsection{9.6 Recommended crystal conventions (practical)}

These are conventions (not hard requirements) that make crystal navigation consistent:

\begin{itemize}
\item **QCM**
\item interference identities → `Σ·✿·INVARIANTS·\{0..2\}`
\item phase lattice objects → `Ψ·✿·OBJECTS·\{0..2\}`
\item bridges (Amp/Meas/Fourier) → `Ω·✿·OPERATORS·\{1..3\}`
\end{itemize}

\begin{itemize}
\item **Gateway**
\item gateway scalars / boosts → `Ω·✿·OPERATORS·0`
\item invariants / calibration identities → `Ω·□·INVARIANTS·\{1..3\}`
\end{itemize}

\begin{itemize}
\item **AQM**
\item Q-number objects → `Ω·☁·OBJECTS·\{0..2\}`
\item meaning transport operators → `Ω·✿·OPERATORS·\{0..2\}`
\item classical shadows / measurement invariants → `Σ·☁·INVARIANTS·\{0..2\}`
\end{itemize}

\begin{itemize}
\item **LM**
\item frames/regimes/state spaces → `Ψ·□·OBJECTS·\{0..2\}`
\item desire/lookahead operators → `Ψ·❋·OPERATORS·\{1..3\}`
\item closure metrics / OICF invariants → `Ψ·☁·INVARIANTS·\{0..2\}`
\end{itemize}

\begin{itemize}
\item **Kernel / Proof**
\item solver outputs → `D·□·CERTIFICATES·\{1..3\}`
\item replay/cert bundles → `D·□·CERTIFICATES·3`
\end{itemize}

The best workflow is:
1) start with “reasonable” assignments,  
2) refine as your atlas grows,  
3) rely on dependency edges + tags to keep everything searchable during refactors.

\_End of manual.\_
\end{document}